\chapter{Le sell side}\label{ch:m1:the-sell-side}

À 11 h 40, un fonds de pension décide de vendre deux millions d'actions d'un titre qui
s'échange à dix millions de titres par jour. Exécuté patiemment, l'ordre prendrait deux jours et
le cours risquerait de baisser pendant ce temps. Le trader du fonds appelle plutôt une banque.
Quarante secondes plus tard, elle obtient un prix pour le bloc entier, un dollar et demi
sous le cours affiché à l'écran, et dit «~conclu~». La banque détient désormais cent millions
de dollars d'actions qu'elle n'a jamais voulues. Ce chapitre traite des sociétés qui
répondent oui à cet appel, de la manière dont elles fixent le prix et des raisons
pour lesquelles elles sont organisées ainsi.

\section{La division marchés}

\begin{definition}[Sell side et broker-dealer]\label{def:m1:the-sell-side:sell-side}
Le \emph{sell side}\index{sell side} est l'ensemble des sociétés qui fournissent des
services de trading aux investisseurs~: exécution, prix, financement, recherche et nouvelles
émissions. Son véhicule juridique est le \emph{broker-dealer}\index{broker-dealer}, une
société agréée à la fois pour exécuter les ordres des clients en qualité d'agent (courtier) et
pour négocier avec eux en qualité de principal (dealer).
\end{definition}

Les plus grands \omterm{def:m1:the-sell-side:sell-side}{broker-dealers} sont logés dans les banques. La \emph{division
marchés} d'une banque est organisée autour du client, et non d'une stratégie~:

\begin{itemize}
\item \textbf{Sales} détient la relation avec le client et apporte la demande.
\item \textbf{Trading} fixe le prix et gère le risque qui en résulte.
\item \textbf{Structuration} conçoit des produits qui regroupent plusieurs risques pour un
  client --- le sujet du livre~5.
\item \textbf{Recherche} publie des analyses qui donnent aux clients une raison de
  les solliciter.
\item \textbf{Quants et technologie} construisent les systèmes de pricing, de gestion des risques et d'exécution
  sur lesquels reposent les quatre autres.
\end{itemize}

\begin{definition}[Sales-trader]\label{def:m1:the-sell-side:sales-trader}
Un \emph{sales-trader}\index{sales-trader} prend les ordres des clients et les travaille
sur le marché pour leur compte, en qualité d'agent, en conseillant sur le moment et les
tactiques~; un \emph{trader}, à l'inverse, engage le capital de la banque.
\end{definition}

\begin{omfigure}
\begin{tikzpicture}[font=\small,
  box/.style={draw=omDef,rounded corners=2pt,minimum width=2.7cm,minimum height=0.85cm,align=center,fill=omDef!4},
  ext/.style={draw=omExo,rounded corners=2pt,minimum width=2.7cm,minimum height=0.85cm,align=center,fill=omExo!6},
  flow/.style={-{Stealth[length=2mm]},thick,omDef}]
\node[ext] (cl) at (0,0) {Client};
\node[box] (sa) at (4.2,1.5) {Sales};
\node[box] (st) at (4.2,0) {Sales-trading\\(en agent)};
\node[box] (re) at (4.2,-1.5) {Recherche};
\node[box] (tr) at (8.6,0.75) {Trading\\(en principal)};
\node[box] (su) at (8.6,-0.75) {Structuration};
\node[ext] (mk) at (13,0.75) {Marchés};
\node[box] (qt) at (8.6,-2.4) {Quants et technologie};
\draw[flow] (cl) -- node[above,sloped]{demande} (sa);
\draw[flow] (cl) -- node[above]{ordres} (st);
\draw[flow] (re) -- node[below,sloped]{idées} (cl);
\draw[flow] (sa) -- (tr);
\draw[flow] (st) -- (tr);
\draw[flow] (tr) -- node[above,font=\footnotesize]{couverture} (mk);
\draw[flow] (su) -- (tr);
\draw[dashed,omExo] (qt) -- (su);
\end{tikzpicture}

\omcaption{Une division marchés vue du client. Tout ce qui figure au centre
sert à apporter de l'activité à la colonne de droite, qui porte le
risque.}
\end{omfigure}

\begin{definition}[Trading de flux et franchise]\label{def:m1:the-sell-side:flow}
Le \emph{trading de flux}\index{trading de flux} est la tenue de marché auprès des clients~: le desk
cote sur demande, prend la contrepartie et gère l'inventaire qui
en résulte. La \emph{franchise}\index{franchise} du desk est le flux des demandes
de clients qu'il reçoit~; sa valeur tient à la fois à la fourchette gagnée sur ce flux
et à l'information que ce flux révèle sur qui veut faire quoi.
\end{definition}

\begin{remark}[Un desk bancaire n'est pas une société de trading pour compte propre]\label{rem:m1:the-sell-side:notprop}
Un \omterm{def:m1:what-a-trading-firm-does:market-maker}{teneur de marché} électronique (\cref{ch:m1:what-a-trading-firm-does}) rivalise
sur le prix et la vitesse pour des ordres anonymes. Un desk bancaire, lui, rivalise pour
\emph{relations}~: le client qui reçoit un bon prix sur un bloc difficile
lui confie aussi à la banque son financement, ses dérivés et sa prochaine
émission obligataire. Une transaction qui perd de l'argent prise isolément peut être une bonne affaire~; décider
quand il faut le faire relève du responsable du desk (livre~16).
\end{remark}

\section{Fixer le prix d'un bloc}

\begin{definition}[Transaction de bloc et risk bid]\label{def:m1:the-sell-side:block}
Une \emph{transaction de bloc}\index{transaction de bloc} est une transaction bien plus grande que la
taille affichée sur le marché, négociée hors du carnet d'ordres. Dans un
\emph{risk bid}, le dealer achète tout le bloc pour compte propre, à une décote
par rapport au cours de marché~; cette décote rémunère le coût et le risque de
la revente ultérieure des actions.
\end{definition}

Le problème du dealer comporte deux volets. Vendre une grande quantité fait baisser le
cours~: c'est un coût d'\emph{impact}. Pendant qu'il vend, le cours bouge pour des
raisons qui lui sont propres~: c'est un \emph{risque}.

\begin{method}[La règle empirique de la racine carrée pour l'impact]\label{met:m1:the-sell-side:sqrt}
Le coût moyen de la négociation de $Q$ actions d'un titre de volatilité
quotidienne $\sigma$ et de volume quotidien $V$, exprimé en fraction de leur valeur, est proche de
\[
  I(Q) \;=\; \kappa\,\sigma\sqrt{Q/V},
\]
avec $\kappa$ de l'ordre de l'unité et, dans des limites raisonnables, peu de dépendance
à la durée de l'exécution. La loi empirique et sa théorie relèvent du
livre~10~; ici, c'est un outil, avec $\kappa = 0.7$.
\end{method}

\begin{proposition}[Risque d'un dénouement linéaire]\label{prop:m1:the-sell-side:tover3}
Un dealer détient un bloc valant $N$ et le vend à un rythme constant sur $T$
jours, tandis que le cours suit une marche aléatoire sans dérive, de volatilité quotidienne
$\sigma$. En négligeant l'impact, son profit par rapport à la vente de l'ensemble au
prix initial a une espérance nulle et un écart-type
\[
  N\,\sigma\sqrt{T/3}.
\]
\end{proposition}
\begin{proof}
Soit $W_t$ la rentabilité du cours depuis l'achat, un mouvement brownien de
variance $\sigma^2 t$. En vendant au rythme $1/T$, le dealer réalise la rentabilité
moyenne $\bar W = \frac1T\int_0^T W_t\dd t$. Alors $\E[\bar W]=0$ et
\[
\Var(\bar W) = \frac{1}{T^2}\int_0^T\!\!\int_0^T \sigma^2 \min(s,t)\dd s\dd t
 = \frac{\sigma^2}{T^2}\cdot\frac{T^3}{3} = \frac{\sigma^2 T}{3}.
\qedhere
\]
\end{proof}

Le facteur $1/3$ est à retenir~: une position qui se réduit linéairement à
zéro sur $T$ jours porte le risque d'une position constante détenue pendant $T/3$
jours.

\begin{proposition}[Décote de seuil de rentabilité]\label{prop:m1:the-sell-side:discount}
Avec un taux de participation $\pi$ (la fraction du volume quotidien que le dealer accepte
de représenter), le dénouement dure $T = Q/(\pi V)$ jours, et la décote pour laquelle le dealer
ne perd de l'argent qu'avec la probabilité $1-\Phi(z)$ seulement est
\[
  d_z \;=\; \kappa\,\sigma\sqrt{Q/V} \;+\; z\,\sigma\sqrt{\frac{Q}{3\pi V}} .
\]
\end{proposition}
\begin{proof}
Le profit du dealer par unité de valeur est $d - I(Q) + \bar W$, une loi normale de
moyenne $d - I(Q)$ et d'écart-type $\sigma\sqrt{T/3}$ d'après la
\cref{prop:m1:the-sell-side:tover3}. Il est négatif avec la probabilité
$\Phi\bigl(-(d-I)/(\sigma\sqrt{T/3})\bigr)$~; on pose cette quantité égale à $\Phi(-z)$.
\end{proof}

\begin{example}[Le bloc du fonds de pension]\label{ex:m1:the-sell-side:block}
$Q = 2$ millions d'actions à 50~\$, $V = 10$ millions, $\sigma = 2\%$ par jour,
$\pi = 10\%$. Alors $T = 2$ jours, $I = 0.7 \times 0.02 \times \sqrt{0.2} =
0.63\%$, et le terme de risque à 95~\% ($z = 1.645$) vaut $1.645 \times 0.02
\times \sqrt{2/3} = 2.69\%$. La décote de seuil de rentabilité est de $3.31\%$, soit 1.66~\$ par
action~: le dealer qui affiche un cours acheteur de 48.34~\$ espère gagner $2.69\%$ sur 100
millions de dollars, soit 2.7~millions de dollars, et perd de l'argent une fois sur vingt. Les deux termes croissent
comme $\sqrt{Q}$~: \emph{un bloc quatre fois plus grand coûte deux fois plus par
action.}
\end{example}

\begin{omfigure}
\begin{tikzpicture}
\begin{axis}[width=11cm,height=5.4cm,
  xlabel={taille du bloc (\% du volume quotidien)},ylabel={décote (\% de la valeur)},
  tick label style={font=\small},label style={font=\small},
  xmin=0,xmax=50,ymin=0,ymax=8.5,grid=major,grid style={gray!25},
  legend columns=4,legend style={font=\footnotesize,at={(0.5,-0.3)},anchor=north,draw=none,fill=none,/tikz/every even column/.append style={column sep=8pt}},
  table/col sep=comma]
\addplot[omThm,very thick] table[x=pct_adv,y=vol3]{figdata/markets-1/02-the-sell-side/discount_vs_size.csv};
\addplot[omDef,very thick] table[x=pct_adv,y=vol2]{figdata/markets-1/02-the-sell-side/discount_vs_size.csv};
\addplot[omMeth,very thick] table[x=pct_adv,y=vol1]{figdata/markets-1/02-the-sell-side/discount_vs_size.csv};
\addplot[omDef,dashed] table[x=pct_adv,y=impact2]{figdata/markets-1/02-the-sell-side/discount_vs_size.csv};
\legend{$\sigma=3\%$,$\sigma=2\%$,$\sigma=1\%$,impact seul ($\sigma=2\%$)}
\end{axis}
\end{tikzpicture}

\omcaption{La décote de seuil à 95~\% d'un cours acheteur intégrant le risque, à un taux de participation de 10~\%.
La majeure partie de la décote rémunère le risque, et non l'impact. Données~: calculées par le
script du chapitre à partir de la \cref{prop:m1:the-sell-side:discount}.}
\end{omfigure}

\begin{remark}[Ce que le modèle laisse de côté]\label{rem:m1:the-sell-side:limits}
Trois éléments, que tout desk évalue à la main. \emph{Information~:} un
client qui vend parce qu'il sait quelque chose mérite une décote plus importante ---
la \omterm{def:m1:what-a-trading-firm-does:adverse-selection}{sélection adverse}, encore. \emph{Couverture~:} le dealer peut vendre des contrats à terme sur indices
en quelques secondes et ne garder que le risque propre à l'action~: une couverture qui supprime
la moitié de la variance, ce qui réduit $\sigma$ de 29~\%. \emph{\omterm{def:m1:the-sell-side:flow}{Franchise}~:} le desk peut
afficher un cours acheteur au-delà de son seuil de rentabilité pour gagner les autres affaires du client.
\end{remark}

\section{Financement, commissions et nouvelles émissions}

Faire des prix est l'une des trois activités de la division.

\begin{definition}[Services de prime brokerage]\label{def:m1:the-sell-side:prime}
Les \emph{services de prime brokerage}\index{services de prime brokerage} désignent l'activité bancaire qui
finance les \omterm{def:m1:what-a-trading-firm-does:hedge-fund}{hedge funds}~: elle leur prête des liquidités en contrepartie de leurs positions longues,
leur prête des titres à vendre à découvert, détient et règle leurs portefeuilles, et est
rémunérée par un spread sur chaque solde. Le \Cref{ch:m1:financing} le décortique.
\end{definition}

\begin{definition}[Prise ferme]\label{def:m1:the-sell-side:underwriting}
Dans la \emph{prise ferme}\index{prise ferme}, une banque achète des actions ou
obligations nouvellement émises auprès de l'émetteur et les revend aux investisseurs, en touchant un pourcentage de
la somme levée et en supportant, pendant quelques heures ou quelques jours, le risque que
l'émission ne se place pas. Avec le \emph{conseil} en fusions-acquisitions, elle forme la
\emph{banque d'affaires}, logée à côté de la division marchés parce que les
nouvelles émissions alimentent les desks de trading.
\end{definition}

\begin{dated}{2026-09}{Une division marchés en ses propres chiffres}\label{dat:m1:the-sell-side:gs}
Goldman Sachs a publié pour 2025 un produit net de 41.5~milliards de dollars dans Global Banking
\& Markets, sur 58.3~milliards de dollars pour l'ensemble de la société. Dans ce périmètre~: taux,
devises et matières premières, 14.5~milliards de dollars, dont 10.3~milliards de dollars
d'intermédiation (fixation des cours) et 4.3~milliards de dollars de financement~; les actions
16.5~milliards de dollars, dont 9.3~milliards d'intermédiation et 7.2~milliards
de financement~; les commissions de banque d'affaires atteignent 9.3~milliards de dollars (conseil 4.7~milliards de dollars,
\omterm{def:m1:the-sell-side:underwriting}{prise ferme} sur actions 1.8~milliards de dollars, \omterm{def:m1:the-sell-side:underwriting}{prise ferme} sur dette 2.8~milliards de dollars).
\end{dated}

\begin{omfigure}
\begin{tikzpicture}
\begin{axis}[width=11.5cm,height=5.2cm,xbar,bar width=7pt,
  xlabel={Produit net 2025 (milliards de \$)},
  ytick=data,yticklabels={{Intermédiation FICC},{Financement FICC},{Intermédiation actions},{Financement actions},{Conseil},{Prise ferme d'actions},{Prise ferme de dette}},
  yticklabel style={font=\small},tick label style={font=\small},label style={font=\small},
  y dir=reverse,xmin=0,xmax=12,enlarge y limits=0.1,grid=major,grid style={gray!25},
  table/col sep=comma]
\addplot[fill=omDef!60,draw=omDef] table[x=usd_bn,y=k]{figdata/markets-1/02-the-sell-side/gs_gbm_2025.csv};
\end{axis}
\end{tikzpicture}

\omcaption{Où une grande division marchés a réalisé son produit en 2025.
Le financement --- prêts adossés aux positions --- représente plus d'un tiers du
produit des activités de trading. Données~: communiqué de résultats de la société,
comme dans l'\cref{dat:m1:the-sell-side:gs}.}
\end{omfigure}

\section{Ce que la réglementation a fait au dealer}

Avant la crise de 2008, les grands dealers négociaient aussi pour leur propre compte sur
des desks dédiés aux opérations pour compte propre. Deux séries de
règles ont mis fin à cette pratique et ont remodelé ce qui subsistait.

\begin{definition}[Actifs pondérés par les risques]\label{def:m1:the-sell-side:rwa}
Les \emph{actifs pondérés par les risques}\index{actifs pondérés par les risques} (RWA, pour risk-weighted assets) d'une banque sont
ses expositions, chacune multipliée par une pondération réglementaire qui reflète son risque~;
la banque doit détenir des fonds propres d'au moins un pourcentage fixe de leur
somme. Pour un desk de trading, la part de RWA liée au risque de marché se calcule à partir du
risque de ses positions selon la norme du Comité de Bâle sur le risque de marché.
\end{definition}

\begin{method}[Juger un desk à l'aune de sa rentabilité des capitaux]\label{met:m1:the-sell-side:roc}
\begin{enumerate}
\item Prendre le produit du desk diminué de ses coûts directs.
\item Diviser par les fonds propres que consomment ses positions~: les RWA multipliés par le
  ratio de fonds propres cible de la banque.
\item Comparer au coût des fonds propres de la banque, d'environ dix pour cent.
\end{enumerate}
Une position qui immobilise des capitaux pendant des semaines doit rapporter bien plus
qu'une position qui tourne en quelques minutes~: c'est pourquoi les desks bancaires
privilégient les flux et le financement plutôt que les risques importants et lents.
\end{method}

\begin{example}[Deux desks]\label{ex:m1:the-sell-side:twodesks}
Le desk A dégage 60~millions de dollars de produit sur des RWA de 2~milliards de dollars~; le desk B dégage 25~millions de dollars
sur des RWA de 400~millions de dollars. Avec un ratio de fonds propres de 12~\%, A mobilise 240~millions de dollars de fonds propres
et rapporte 25~\%~; B en mobilise 48~millions de dollars et rapporte 52~\%. La banque développe B.
\end{example}

La seconde règle est une interdiction. La section~13 du Bank Holding Company Act américain,
ajoutée par le Dodd--Frank Act de 2010 et connue sous le nom de \emph{Volcker
rule}, interdit aux entités bancaires de se livrer en principal à des opérations de
trading pour compte propre à court terme. Ses exceptions définissent le dealer moderne~: la \omterm{def:m1:the-sell-side:underwriting}{prise ferme},
l'activité liée au \emph{market making}, la couverture de réduction des risques, le trading
pour le compte de clients et le trading d'obligations d'État restent autorisés.
Un desk doit donc pouvoir démontrer que son stock est dimensionné pour répondre aux
«~demandes raisonnablement prévisibles à court terme des clients, des consommateurs ou des
contreparties~» --- autrement dit qu'il est une boutique et non un fonds. Une grande partie des talents issus du trading pour compte propre qui ont quitté les banques après 2010
ont fondé ou rejoint les sociétés décrites au \cref{ch:m1:what-a-trading-firm-does}.

\section{Tutoriel~: tarifer le risk bid}\label{tut:m1:the-sell-side:riskbid}

\begin{tutorial}
\textbf{Objectif.} Tarifer le bloc du fonds de pension et vérifier la règle du $T/3$ par
simulation. \textbf{État final~:} l'histogramme ci-dessous, avec 5~\% de sa masse
à gauche de zéro.
\begin{tutsteps}
\item \textbf{Coder les deux propositions.}
\omcode{code/markets-1/02-the-sell-side/python/blockbid.py}{29}{45}{Durée de liquidation, impact, risque et décote au seuil de rentabilité.}\label{lst:m1:the-sell-side:formulas}
\item \textbf{Simuler la liquidation.} Deux cents tranches égales sur $T$ jours,
chacune vendue au prix du moment~; le résultat du dealer est la décote, diminuée de
l'impact, augmentée du rendement moyen des tranches.
\omcode{code/markets-1/02-the-sell-side/python/blockbid.py}{47}{54}{Monte Carlo d'une liquidation linéaire~: la moyenne d'une marche aléatoire.}\label{lst:m1:the-sell-side:mc}
\item \textbf{Exécuter et comparer.} Pour le bloc de l'exemple, l'écart type simulé
vaut $1.64\%$, contre $0.02\sqrt{2/3} =
1.63\%$ prédit, et $4.9\%$ des 20\,000 trajectoires perdent de l'argent à la décote de 95~\%.
\item \textbf{Lire l'histogramme.} La moyenne, $2.7\%$, est le profit espéré du
dealer~; c'est ce que le client paie pour la certitude.
\end{tutsteps}
\textbf{Pour aller plus loin.} Augmenter \texttt{participation} à 0.25~: la décote tombe à
$2.3\%$, et le modèle affirme qu'aller plus vite est toujours préférable. Pourquoi aucun desk
ne vend-il à 50~\% du volume~? (Exercice~\ref{exo:m1:the-sell-side:8}.)
Puis couvrir~: remplacer $\sigma$ par la volatilité résiduelle du titre.
\end{tutorial}

\begin{omfigure}
\begin{tikzpicture}
\begin{axis}[width=11cm,height=4.8cm,ybar,bar width=6.5pt,
  xlabel={profit du dealer (\% de la valeur du bloc)},ylabel={trajectoires sur 20\,000},
  tick label style={font=\small},label style={font=\small},
  xmin=-4.2,xmax=9.2,ymin=0,table/col sep=comma]
\addplot[fill=omDef!40,draw=omDef] table[x=centre,y=count]{figdata/markets-1/02-the-sell-side/pnl_hist.csv};
\addplot[omThm,dashed,sharp plot,no markers] coordinates {(0,0) (0,2500)};
\end{axis}
\end{tikzpicture}

\omcaption{Profit du risk bid au seuil de rentabilité de 95~\%. La ligne
pointillée marque zéro~: une trajectoire sur vingt tombe à sa gauche. Données~: la
simulation du tutoriel, graine 2.}
\end{omfigure}

\section{Construction~: le tarificateur de bloc}\label{bld:m1:the-sell-side:blockbid}

\begin{build}
\bfield{Rôle} Un service de tarification que le desk de type banque de la petite société
appelle lorsqu'un client demande un risk bid. \bfield{Interface} \texttt{quote(shares, price, adv, sigma, participation,
confidence, hedge\_ratio=0.0)} renvoyant la durée de liquidation, l'impact, le risque,
la décote et le cours acheteur. \texttt{hedge\_ratio} est la fraction de variance
éliminée par une couverture par indice~: le terme de risque utilise $\sigma\sqrt{1-h}$.
\bfield{Règles} Rejeter une participation supérieure à 30~\% et un bloc supérieur à cinq
jours de volume (\texttt{ValueError})~: hors de cette plage, la règle de la racine carrée
n'est pas fiable. Arrondir le cours acheteur \emph{à la baisse} au centime.
\bfield{Tests d'acceptation} \texttt{code/firm/blockbid/tests/}~: reproduit
l'\cref{ex:m1:the-sell-side:block} au centime près~; un bloc entièrement couvert est tarifé
au seul impact~; les limites lèvent une exception.
\bfield{Pour aller plus loin} Inscrire le profit réalisé par le desk sur une liquidation simulée
dans le registre présenté au \cref{ch:m1:what-a-trading-firm-does} sous \texttt{position}.
\end{build}

\begin{omsources}
\item Goldman Sachs, \emph{Full Year and Fourth Quarter 2025 Earnings
Results}, Form 8-K, janvier 2026.
\item Federal Deposit Insurance Corporation, «~Volcker Rule~» (aperçu de
Section 13 du Bank Holding Company Act et ses règles d'application);
\emph{Federal Register}, 31 juillet 2020, règle finale sur les covered funds.
\item Basel Committee on Banking Supervision, \emph{Minimum capital
requirements for market risk}, janvier 2019.
\item L.~Harris, \emph{Trading and Exchanges}, Oxford University Press, 2003,
chapitres 15 et 16 (transactions de bloc, dealers).
\item J.-P.~Bouchaud, J.~Bonart, J.~Donier et M.~Gould, \emph{Trades, Quotes
and Prices}, Cambridge University Press, 2018, chapitre 12 (la loi de la
racine carrée).
\end{omsources}

\section{Exercices}

\begin{exercise}[$\star$]\label{exo:m1:the-sell-side:1}
Un dealer détient une position de 30~millions de dollars dans un titre dont la volatilité quotidienne
est de 1.5~\%, et il la liquide linéairement sur trois jours. Calculer l'écart type de son
résultat en dollars.
\end{exercise}

\begin{exercise}[$\star$]\label{exo:m1:the-sell-side:2}
Un bloc de 500\,000 actions d'un titre qui s'échange pour 5~millions par jour est liquidé
à une participation de 10~\%. Calculer la durée de la liquidation et le coût d'impact en
pourcentage si $\sigma = 1.5\%$ et $\kappa = 0.7$~?
\end{exercise}

\begin{exercise}[$\star$]\label{exo:m1:the-sell-side:3}
Un émetteur lève 800~millions de dollars d'actions nouvelles, avec une commission de \omterm{def:m1:the-sell-side:underwriting}{prise ferme}
de 3.5~\%. Que gagnent les banques~? Les actions sont placées à 40~dollars~; le premier jour,
elles clôturent à 46~dollars. Quel montant l'émetteur a-t-il «~laissé sur la table~», et
qui l'a reçu~?
\end{exercise}

\begin{exercise}[$\star\star$]\label{exo:m1:the-sell-side:4}
Calculer la décote de seuil de rentabilité à 95~\%, en pourcentage et en dollars par action,
pour un bloc de 3~millions d'actions à 20~dollars dans un titre avec $V = 6$ millions et
$\sigma = 2.5\%$, à une participation de 10~\%.
\end{exercise}

\begin{exercise}[$\star\star$]\label{exo:m1:the-sell-side:5}
Montrer qu'à participation fixée, la décote de seuil est proportionnelle à
$\sqrt{Q}$. Un dealer cote 2~\% pour un bloc~; que devrait-il coter, toutes choses
égales par ailleurs, pour un bloc neuf fois plus grand~?
\end{exercise}

\begin{exercise}[$\star\star$]\label{exo:m1:the-sell-side:6}
À partir des chiffres de l'\cref{dat:m1:the-sell-side:gs}~: (a) quelle part du chiffre d'affaires
de la division provient du financement, de l'intermédiation et des
commissions de banque d'affaires~? (b) À laquelle des cinq sources de revenus de
la \cref{met:m1:what-a-trading-firm-does:sources} correspond chacune~?
\end{exercise}

\begin{exercise}[$\star\star\star$]\label{exo:m1:the-sell-side:7}
\emph{Programmation.} Une couverture par indice supprime 40~\% de la variance du titre de
l'\cref{ex:m1:the-sell-side:block}. (a) Calculer la nouvelle décote de seuil de rentabilité.
(b) La couverture coûte 2~points de base de la valeur du bloc en contrats à terme,
en commissions et en spread. Est-elle rentable~?
\end{exercise}

\begin{exercise}[$\star\star\star$]\label{exo:m1:the-sell-side:8}
\emph{Trouver le défaut.} Un jeune trader constate que dans la
\cref{prop:m1:the-sell-side:discount}, la décote diminue lorsque la participation augmente,
et propose de liquider chaque bloc à 50~\% du volume. Donner deux raisons, l'une à l'intérieur
des hypothèses du modèle et l'autre en dehors, pour lesquelles la décote du desk ne diminuerait pas
comme l'indique la formule.
\end{exercise}

\section{Problème~: le cours acheteur de risque}

\begin{problem}[{Problème du week-end --- quarante secondes pour valoriser un bloc}]\label{pb:m1:the-sell-side:1}
On dirige le livre des actions au comptant et du risque d'une banque. À 11~h~40, un fonds de pension demande
un cours acheteur sur 4~millions d'actions d'un titre qui s'échange à 25~dollars, avec un volume quotidien
de 8~millions d'actions et une volatilité quotidienne de 1.8~\%. Prendre $\kappa = 0.7$.

\medskip\noindent\textbf{Partie I --- Le prix mécanique.}
\begin{enumerate}
\item Quelle est la valeur du bloc, et quelle fraction du volume quotidien représente-t-elle~?
\item À une participation de 10~\%, combien de jours dure la liquidation~?
\item Calculer le coût d'impact en pourcentage et en dollars.
\item Calculer l'écart type du résultat de la liquidation en pourcentage et en
dollars.
\item Calculer la décote de seuil de rentabilité à 95~\% et le cours acheteur correspondant,
arrondi au cent inférieur.
\end{enumerate}

\medskip\noindent\textbf{Partie II --- Améliorer le prix.}
\begin{enumerate}[resume]
\item On accepte d'être 20~\% du volume. Recalculer la durée de liquidation
et la décote.
\item Une couverture par indice supprime la moitié de la variance. Recalculer la décote à
une participation de 20~\%.
\item Quel est le profit espéré en dollars si l'on remporte le bloc à
cette décote~?
\item On sait qu'un concurrent cote à un écart type ($z = 1$)
au lieu de 1.645. Quelle décote propose-t-il avec la même couverture et la même
participation, et à quelle fréquence perd-il de l'argent~?
\end{enumerate}

\medskip\noindent\textbf{Partie III --- Le client.}
\begin{enumerate}[resume]
\item L'alternative pour le fonds est une exécution en agence sur les mêmes
$T$ jours, avec une commission de 1~cent par action. Donner son coût espéré en
pourcentage (impact plus commission) et l'écart type de son résultat.
\item De combien le cours acheteur de risque de la question~7 dépasse-t-il ce
coût espéré~? Qu'achète le fonds avec cette différence~?
\item Le trader du fonds est évalué par rapport au prix au moment de la
décision. Expliquer pourquoi elle peut préférer le cours acheteur de risque, même avec cette prime.
\item Au cours de l'année écoulée, les blocs de ce client ont été suivis, en moyenne, par
une baisse de 0.4~\% du titre par rapport à l'indice dans les deux jours suivants.
Que faut-il ajouter à la décote, et comment appelle-t-on ce coût~?
\end{enumerate}

\medskip\noindent\textbf{Partie IV --- La banque.}
\begin{enumerate}[resume]
\item La position ajoute 150~millions de dollars d'\omterm{def:m1:the-sell-side:rwa}{actifs pondérés par les
risques} pendant trois jours. Avec un ratio de fonds propres de 12~\% et un
coût des fonds propres de 10~\%, calculer le coût de ce capital en dollars (on prendra 252 jours).
\item Comparer avec le profit espéré de la question~8. Le capital est-il
le coût contraignant de cette transaction~?
\item Le même client a versé à la banque 3~millions de dollars l'année dernière au titre
du financement et des produits dérivés. Déterminer le cours acheteur le plus bas que l'on puisse rationnellement proposer.
\item Vous remportez le bloc et vendez contre celui-ci 100~millions de dollars
de contrats à terme sur indice. Dans la première heure, l'indice baisse de 1~\% et le
titre de 3~\%. Estimer la perte en valorisation au prix de marché sur la paire avant toute vente.
\item Expliquer, au regard de la Volcker rule, ce qu'il faudrait démontrer
à un examinateur au sujet de cette position.
\item Énoncer le \emph{résultat nommé}~: la décote, en cents par action, à partir
de laquelle ce desk atteint le seuil de rentabilité avec un niveau de confiance de
95~\%, après couverture et la charge de \omterm{def:m1:what-a-trading-firm-does:adverse-selection}{sélection adverse} de la question~13.
\item En une phrase, expliquer pourquoi un cours acheteur de risque est une option que détient le client.
\end{enumerate}
\end{problem}

\section{Questions d'entretien}

\begin{interviewq}[$\star$ \iqroles{banque, trader}]\label{iq:m1:the-sell-side:1}
Quelle est la différence entre un \omterm{def:m1:the-sell-side:sales-trader}{sales-trader} et un trader~? Qui est
responsable du P\&L d'un ordre client dans chaque cas~?
\end{interviewq}

\begin{interviewq}[$\star\star$ \iqroles{trader, chercheur, banque}]\label{iq:m1:the-sell-side:2}
Vous devez vendre une position de façon régulière sur neuf jours. Environ quel
risque portez-vous, comparé à une détention inchangée pendant les mêmes neuf jours~?
\end{interviewq}

\begin{interviewq}[$\star\star$ \iqroles{trader, banque}]\label{iq:m1:the-sell-side:3}
Un client vous demande une cotation sur un bloc d'un titre que vous ne connaissez pas du tout.
Que vous demandez-vous pendant les dix secondes dont vous disposez~?
\end{interviewq}

\begin{interviewq}[$\star\star$ \iqroles{banque, chercheur}]\label{iq:m1:the-sell-side:4}
Pourquoi le financement a-t-il représenté plus d'un tiers du revenu de négoce
des grandes banques après 2010, alors que leur risque pour compte propre diminuait~?
\end{interviewq}

\begin{interviewq}[$\star\star\star$ \iqroles{chercheur, trader}]\label{iq:m1:the-sell-side:5}
Le coût d'un bloc croît comme la racine carrée de sa taille. Un client peut
vous confier un bloc de 4~millions d'actions ou quatre blocs de 1~million sur
quatre jours consécutifs. Lequel est le moins cher pour le client, et quelle
hypothèse le décide~?
\end{interviewq}

\begin{interviewq}[$\star\star\star$ \iqroles{banque, chercheur}]\label{iq:m1:the-sell-side:6}
Calculez la variance de la moyenne temporelle d'un mouvement brownien sur $[0,T]$.
\end{interviewq}
